% !TEX program = xelatex
\def\StudyRoot{../..}
\def\StudyBook{WrittenResponse}
\documentclass[11pt,lang=cn,chinesefont=nofont,mode=simple,math=newtx,color=blue]{elegantbook}
\usepackage{longtable,array,fvextra,xurl,adjustbox,needspace}
\usepackage[normalem]{ulem}
\providecommand{\lt}{<}
\providecommand{\gt}{>}
\tracinglostchars=3
\setCJKmainfont{FandolSong-Regular.otf}[BoldFont=FandolSong-Bold.otf,ItalicFont=FandolKai-Regular.otf]
\setCJKsansfont{FandolHei-Regular.otf}[BoldFont=FandolHei-Bold.otf]
\setCJKmonofont{FandolFang-Regular.otf}
\setmonofont{DejaVuSansMono.ttf}[BoldFont=DejaVuSansMono-Bold.ttf,ItalicFont=DejaVuSansMono-Oblique.ttf,Scale=.84]
\geometry{a4paper,left=22mm,right=18mm,top=20mm,bottom=21mm,headheight=15pt,headsep=7mm}
\definecolor{structurecolor}{HTML}{173F57}
\colorlet{main}{structurecolor}
\colorlet{second}{structurecolor}
\colorlet{third}{structurecolor}
\hypersetup{unicode=true,colorlinks=true,linkcolor=structurecolor,urlcolor=structurecolor,citecolor=structurecolor,pdfauthor={CityU StudyNotes}}
\makeatletter
\renewcommand{\tableofcontents}{%
  \hypersetup{linktoc=all,linkcolor=black}%
  \chapter*{\contentsname}%
  \StudyTarget{\StudyBook__contents}\StudyTarget{\StudyBook___1}%
  \@starttoc{toc}%
  \hypersetup{linkcolor=structurecolor}\clearpage}
\makeatother
\setstretch{1.25}
\setlength{\parskip}{.35\baselineskip}
\setlength{\parindent}{0pt}
\setlength{\emergencystretch}{2em}
\setcounter{tocdepth}{1}
\setcounter{secnumdepth}{-1}
\setlist{nosep,leftmargin=1.8em,topsep=.35\baselineskip}
\titleformat{\chapter}[display]{\Huge\sffamily\bfseries\color{structurecolor}}{}{0pt}{}[\vspace{.4ex}\titlerule]
\titlespacing*{\chapter}{0pt}{0pt}{22pt}
\titleformat{\section}{\Large\sffamily\bfseries\color{structurecolor}}{}{0pt}{}
\titleformat{\subsection}{\large\sffamily\bfseries}{}{0pt}{}
\titleformat{\subsubsection}{\normalsize\sffamily\bfseries}{}{0pt}{}
\titlespacing*{\section}{0pt}{1.6\baselineskip}{.5\baselineskip}
\titlespacing*{\subsection}{0pt}{1.1\baselineskip}{.35\baselineskip}
\providecommand{\StudyBook}{Sample}
\providecommand{\StudyRoot}{../..}
\providecommand{\StudySourceRef}{main}
\providecommand{\StudyDate}{2026-10-16}
\newcommand{\StudyRunning}{\StudyBook}
\fancyhf{}
\fancyhead[L]{\small\sffamily\color{gray}\StudyRunning}
\fancyhead[R]{\small\sffamily\color{gray}CityU StudyNotes}
\fancyfoot[R]{\small\thepage}
\renewcommand{\headrulewidth}{.3pt}
\pagestyle{fancy}
\fancypagestyle{plain}{\fancyhf{}\fancyfoot[R]{\small\thepage}\renewcommand{\headrulewidth}{0pt}}
\newwrite\StudyMap
\AtBeginDocument{\immediate\openout\StudyMap=\jobname.pages}
\makeatletter
\newcommand{\StudyTarget}[1]{%
  \hypertarget{#1}{}\label{#1}%
  \protected@write\StudyMap{}{#1|\thepage}%
}
\makeatother
\newcommand{\StudyChapter}[2]{%
  \clearpage\renewcommand{\StudyRunning}{\ifcsname studytitle@#1\endcsname\csname studytitle@#1\endcsname\else#1\fi}%
  \chapter*{#2}\phantomsection\addcontentsline{toc}{chapter}{#2}%
  \StudyTarget{#1__root}%
}
\newcommand{\StudyHeading}[2]{%
  \ifcase#1\relax\or
    \section*{#2}\phantomsection\addcontentsline{toc}{section}{#2}%
  \or\subsection*{#2}\or\subsubsection*{#2}\else\paragraph*{#2}\fi
}
\newcommand{\StudyImage}[1]{%
  \par\medskip\begingroup
  \edef\StudyResolvedImage{\ifcsname studyimage@#1\endcsname\csname studyimage@#1\endcsname\else\StudyRoot/#1\fi}%
  \noindent\makebox[\linewidth][c]{%
  \adjustbox{max width=.96\linewidth,max height=.69\textheight}{%
  \includegraphics{\StudyResolvedImage}}}\endgroup\par\medskip
}
\newcommand{\StudyCellImage}[1]{%
  \begingroup
  \edef\StudyResolvedImage{\ifcsname studyimage@#1\endcsname\csname studyimage@#1\endcsname\else\StudyRoot/#1\fi}%
  \adjustbox{max width=\linewidth,max height=.28\textheight}{\includegraphics{\StudyResolvedImage}}\endgroup}
\DefineVerbatimEnvironment{StudyCode}{Verbatim}{breaklines=true,breakanywhere=true,fontsize=\small,tabsize=4,framesep=3mm,frame=leftline,rulecolor=\color{gray!55}}
\newcommand{\StudyInline}[1]{\allowbreak{}\texttt{#1}\allowbreak{}}
\newenvironment{StudySolution}[1]{%
  \par\medskip\Needspace{3\baselineskip}\noindent
  {\bfseries\color{structurecolor}#1}\par\nobreak\smallskip
}{\par\medskip}
\newcommand{\StudyBadge}[2]{%
  \begingroup\small\sffamily
  \ifstrequal{#1}{mid}{\colorlet{studybadge}{green!35!black}}{%
  \ifstrequal{#1}{final}{\colorlet{studybadge}{yellow!45!black}}{\colorlet{studybadge}{red!45!black}}}%
  \textcolor{studybadge}{\rule[-.2ex]{1.2pt}{1.2em}}\enspace\textcolor{black}{#2}\quad\endgroup}
\newcommand{\StudyRef}[3]{%
  \ifcsname studyref@#1@#2\endcsname
    \csname studyref@#1@#2\endcsname{#3}%
  \else #3\textsuperscript{[待定页]}\fi}
\newcommand{\StudyRepoRef}[3]{%
  \href{https://github.com/CrazyShout/CityU-StudyNotes/blob/\StudySourceRef/#1\ifstrempty{#2}{}{\##2}}{#3}}
\newcommand{\StudyNotebookCell}[2]{%
  \ifcsname studynb@#1@#2\endcsname
    \input{\csname studynb@#1@#2\endcsname}%
  \else\PackageError{studynotes}{Missing saved notebook cell #2}{}\fi}
\newcommand{\StudyCoverExtra}{}
\newcommand{\StudyCover}[2]{%
  \thispagestyle{empty}
  \StudyTarget{\StudyBook__root}\StudyCoverExtra
  \vspace*{22mm}
  {\sffamily\large\color{structurecolor}CITYU STUDYNOTES\par}
  \vspace{12mm}
  {\sffamily\bfseries\Huge\color{structurecolor}#1\par}
  \vspace{6mm}{\Large #2\par}
  \vspace{10mm}\noindent\textcolor{structurecolor}{\rule{.42\linewidth}{1pt}}\par
  \vspace{8mm}中文讲解 · English terminology · 双语题答\par
  A4 单栏 · 根据已下载教师材料整理\par
  \vfill
  {\small 本资料用于课程学习与复习，保留原材料、历史题源与教学补充的身份。\par
  模板：ElegantBook 4.7；版式与构建由 CityU-StudyNotes 维护。\par
  \StudyDate\par}
  \clearpage
}
\newcommand{\StudyQuestionExtra}{}
\newcommand{\StudyQuestion}[3]{%
  \Needspace{10\baselineskip}
  \section*{#1 · #2}\phantomsection\addcontentsline{toc}{section}{#1 · #2}%
  \StudyTarget{\StudyBook__#1}%
  \StudyQuestionExtra
  \lowercase{\def\StudyQuestionAnchor{#1}}%
  \noindent\begin{minipage}[t]{.65\linewidth}{\small\sffamily #3}\end{minipage}\hfill
  \begin{minipage}[t]{.33\linewidth}\raggedleft\scriptsize
  \StudyRef{Questions}{\StudyQuestionAnchor}{题目} · \StudyRef{Answers}{\StudyQuestionAnchor}{答案}
  \end{minipage}\par\medskip
}
\InputIfFileExists{\StudyRoot/.build/latex/\StudyBook/refs.tex}{}{}

\begin{document}
\setlength{\parskip}{.25\baselineskip}
\titlespacing*{\subsection}{0pt}{.7\baselineskip}{.2\baselineskip}
\renewcommand{\StudyCoverExtra}{\StudyTarget{WrittenResponse__cs5489-written-response-questions-and-answers}}
\StudyCover{CS5489 · 期中简答题与答案}{Written-Response Questions and Answers}
\input{\StudyRoot/CS5489/exam-review/Midterm/frontmatter/WrittenResponse.tex}
\begingroup\small\setlength{\cftbeforechapskip}{.45\baselineskip}\tableofcontents\endgroup
\Needspace{7\baselineskip}\phantomsection\addcontentsline{toc}{chapter}{学习流程与模型诊断}
\renewcommand{\StudyQuestionExtra}{\StudyTarget{WrittenResponse__mt011}}
\StudyQuestion{MT011}{识别过拟合，并提出可操作的改进}{简答题 / Written response · 解释、提出办法}
\input{\StudyRoot/CS5489/exam-review/Midterm/items/MT011/provenance.tex}
\subsection*{English question}
\StudyTarget{WrittenResponse__english-question}
\input{\StudyRoot/CS5489/exam-review/Midterm/items/MT011/question-en.tex}
\subsection*{中文题意}
\StudyTarget{WrittenResponse___2}
\input{\StudyRoot/CS5489/exam-review/Midterm/items/MT011/question-zh.tex}
\subsection*{English answer}
\StudyTarget{WrittenResponse__english-answer}
\input{\StudyRoot/CS5489/exam-review/Midterm/items/MT011/answer-en.tex}
\subsection*{中文对应答案}
\StudyTarget{WrittenResponse___3}
\input{\StudyRoot/CS5489/exam-review/Midterm/items/MT011/answer-zh.tex}
\input{\StudyRoot/CS5489/exam-review/Midterm/items/MT011/concise-extra.tex}
\renewcommand{\StudyQuestionExtra}{\StudyTarget{WrittenResponse__mt020}}
\StudyQuestion{MT020}{数据很多，训练和测试仍都差}{简答题 / Written response · 诊断问题、提出办法}
\input{\StudyRoot/CS5489/exam-review/Midterm/items/MT020/provenance.tex}
\subsection*{English question}
\StudyTarget{WrittenResponse__english-question_1}
\input{\StudyRoot/CS5489/exam-review/Midterm/items/MT020/question-en.tex}
\subsection*{中文题意}
\StudyTarget{WrittenResponse___4}
\input{\StudyRoot/CS5489/exam-review/Midterm/items/MT020/question-zh.tex}
\subsection*{English answer}
\StudyTarget{WrittenResponse__english-answer_1}
\input{\StudyRoot/CS5489/exam-review/Midterm/items/MT020/answer-en.tex}
\subsection*{中文对应答案}
\StudyTarget{WrittenResponse___5}
\input{\StudyRoot/CS5489/exam-review/Midterm/items/MT020/answer-zh.tex}
\input{\StudyRoot/CS5489/exam-review/Midterm/items/MT020/concise-extra.tex}
\renewcommand{\StudyQuestionExtra}{\StudyTarget{WrittenResponse__mt022}}
\StudyQuestion{MT022}{五万客户、需要解释：选择哪种分类器？}{简答题 / Written response · 比较、提出办法}
\input{\StudyRoot/CS5489/exam-review/Midterm/items/MT022/provenance.tex}
\subsection*{English question}
\StudyTarget{WrittenResponse__english-question_2}
\input{\StudyRoot/CS5489/exam-review/Midterm/items/MT022/question-en.tex}
\subsection*{中文题意}
\StudyTarget{WrittenResponse___6}
\input{\StudyRoot/CS5489/exam-review/Midterm/items/MT022/question-zh.tex}
\subsection*{English answer}
\StudyTarget{WrittenResponse__english-answer_2}
\input{\StudyRoot/CS5489/exam-review/Midterm/items/MT022/answer-en.tex}
\subsection*{中文对应答案}
\StudyTarget{WrittenResponse___7}
\input{\StudyRoot/CS5489/exam-review/Midterm/items/MT022/answer-zh.tex}
\input{\StudyRoot/CS5489/exam-review/Midterm/items/MT022/concise-extra.tex}
\Needspace{7\baselineskip}\phantomsection\addcontentsline{toc}{chapter}{Bayes、Naive Bayes与模型比较}
\renewcommand{\StudyQuestionExtra}{\StudyTarget{WrittenResponse__mt046}\StudyTarget{WrittenResponse__l2q09}\StudyTarget{WrittenResponse__mt046-gaussian-nb}}
\StudyQuestion{MT046}{完整讲出Gaussian NB}{简答题 / Written response · 解释}
\input{\StudyRoot/CS5489/exam-review/Midterm/items/MT046/provenance.tex}
\subsection*{English question}
\StudyTarget{WrittenResponse__english-question_3}
\input{\StudyRoot/CS5489/exam-review/Midterm/items/MT046/question-en.tex}
\subsection*{中文题意}
\StudyTarget{WrittenResponse___8}
\input{\StudyRoot/CS5489/exam-review/Midterm/items/MT046/question-zh.tex}
\subsection*{English answer}
\StudyTarget{WrittenResponse__english-answer_3}
\input{\StudyRoot/CS5489/exam-review/Midterm/items/MT046/answer-en.tex}
\subsection*{中文对应答案}
\StudyTarget{WrittenResponse___9}
\input{\StudyRoot/CS5489/exam-review/Midterm/items/MT046/answer-zh.tex}
\input{\StudyRoot/CS5489/exam-review/Midterm/items/MT046/concise-extra.tex}
\renewcommand{\StudyQuestionExtra}{\StudyTarget{WrittenResponse__mt034}\StudyTarget{WrittenResponse__l2q15}}
\StudyQuestion{MT034}{无限数据能消除分类错误吗}{简答题 / Written response · 判断并说明}
\input{\StudyRoot/CS5489/exam-review/Midterm/items/MT034/provenance.tex}
\subsection*{English question}
\StudyTarget{WrittenResponse__english-question_4}
\input{\StudyRoot/CS5489/exam-review/Midterm/items/MT034/question-en.tex}
\subsection*{中文题意}
\StudyTarget{WrittenResponse___10}
\input{\StudyRoot/CS5489/exam-review/Midterm/items/MT034/question-zh.tex}
\subsection*{English answer}
\StudyTarget{WrittenResponse__english-answer_4}
\input{\StudyRoot/CS5489/exam-review/Midterm/items/MT034/answer-en.tex}
\subsection*{中文对应答案}
\StudyTarget{WrittenResponse___11}
\input{\StudyRoot/CS5489/exam-review/Midterm/items/MT034/answer-zh.tex}
\input{\StudyRoot/CS5489/exam-review/Midterm/items/MT034/concise-extra.tex}
\renewcommand{\StudyQuestionExtra}{\StudyTarget{WrittenResponse__mt033}\StudyTarget{WrittenResponse__l2q04}}
\StudyQuestion{MT033}{用图像特征解释两种建模方式}{简答题 / Written response · 比较、解释}
\input{\StudyRoot/CS5489/exam-review/Midterm/items/MT033/provenance.tex}
\subsection*{English question}
\StudyTarget{WrittenResponse__english-question_5}
\input{\StudyRoot/CS5489/exam-review/Midterm/items/MT033/question-en.tex}
\subsection*{中文题意}
\StudyTarget{WrittenResponse___12}
\input{\StudyRoot/CS5489/exam-review/Midterm/items/MT033/question-zh.tex}
\subsection*{English answer}
\StudyTarget{WrittenResponse__english-answer_5}
\input{\StudyRoot/CS5489/exam-review/Midterm/items/MT033/answer-en.tex}
\subsection*{中文对应答案}
\StudyTarget{WrittenResponse___13}
\input{\StudyRoot/CS5489/exam-review/Midterm/items/MT033/answer-zh.tex}
\input{\StudyRoot/CS5489/exam-review/Midterm/items/MT033/concise-extra.tex}
\renewcommand{\StudyQuestionExtra}{\StudyTarget{WrittenResponse__mt059}\StudyTarget{WrittenResponse__l2q13}\StudyTarget{WrittenResponse__mt059-lrgaussian-nb}}
\StudyQuestion{MT059}{LR与共享方差Gaussian NB}{简答题 / Written response · 比较}
\input{\StudyRoot/CS5489/exam-review/Midterm/items/MT059/provenance.tex}
\subsection*{English question}
\StudyTarget{WrittenResponse__english-question_6}
\input{\StudyRoot/CS5489/exam-review/Midterm/items/MT059/question-en.tex}
\subsection*{中文题意}
\StudyTarget{WrittenResponse___14}
\input{\StudyRoot/CS5489/exam-review/Midterm/items/MT059/question-zh.tex}
\subsection*{English answer}
\StudyTarget{WrittenResponse__english-answer_6}
\input{\StudyRoot/CS5489/exam-review/Midterm/items/MT059/answer-en.tex}
\subsection*{中文对应答案}
\StudyTarget{WrittenResponse___15}
\input{\StudyRoot/CS5489/exam-review/Midterm/items/MT059/answer-zh.tex}
\input{\StudyRoot/CS5489/exam-review/Midterm/items/MT059/concise-extra.tex}
\renewcommand{\StudyQuestionExtra}{\StudyTarget{WrittenResponse__mt021}\StudyTarget{WrittenResponse__l2q12}}
\StudyQuestion{MT021}{同为线性分类器，学习原则有何不同}{简答题 / Written response · 比较}
\input{\StudyRoot/CS5489/exam-review/Midterm/items/MT021/provenance.tex}
\subsection*{English question}
\StudyTarget{WrittenResponse__english-question_7}
\input{\StudyRoot/CS5489/exam-review/Midterm/items/MT021/question-en.tex}
\subsection*{中文题意}
\StudyTarget{WrittenResponse___16}
\input{\StudyRoot/CS5489/exam-review/Midterm/items/MT021/question-zh.tex}
\subsection*{English answer}
\StudyTarget{WrittenResponse__english-answer_7}
\input{\StudyRoot/CS5489/exam-review/Midterm/items/MT021/answer-en.tex}
\subsection*{中文对应答案}
\StudyTarget{WrittenResponse___17}
\input{\StudyRoot/CS5489/exam-review/Midterm/items/MT021/answer-zh.tex}
\input{\StudyRoot/CS5489/exam-review/Midterm/items/MT021/concise-extra.tex}
\renewcommand{\StudyQuestionExtra}{\StudyTarget{WrittenResponse__mt058}\StudyTarget{WrittenResponse__l2q16}\StudyTarget{WrittenResponse__mt058-nb}}
\StudyQuestion{MT058}{智能笔上的文本NB如何改进}{简答题 / Written response · 提出办法}
\input{\StudyRoot/CS5489/exam-review/Midterm/items/MT058/provenance.tex}
\subsection*{English question}
\StudyTarget{WrittenResponse__english-question_8}
\input{\StudyRoot/CS5489/exam-review/Midterm/items/MT058/question-en.tex}
\subsection*{中文题意}
\StudyTarget{WrittenResponse___18}
\input{\StudyRoot/CS5489/exam-review/Midterm/items/MT058/question-zh.tex}
\subsection*{English answer}
\StudyTarget{WrittenResponse__english-answer_8}
\input{\StudyRoot/CS5489/exam-review/Midterm/items/MT058/answer-en.tex}
\subsection*{中文对应答案}
\StudyTarget{WrittenResponse___19}
\input{\StudyRoot/CS5489/exam-review/Midterm/items/MT058/answer-zh.tex}
\input{\StudyRoot/CS5489/exam-review/Midterm/items/MT058/concise-extra.tex}
\Needspace{7\baselineskip}\phantomsection\addcontentsline{toc}{chapter}{逻辑回归与参数正则}
\renewcommand{\StudyQuestionExtra}{\StudyTarget{WrittenResponse__mt047}\StudyTarget{WrittenResponse__mt047-logistic}}
\StudyQuestion{MT047}{解释Logistic目标中的两项}{简答题 / Written response · 解释}
\input{\StudyRoot/CS5489/exam-review/Midterm/items/MT047/provenance.tex}
\subsection*{English question}
\StudyTarget{WrittenResponse__english-question_9}
\input{\StudyRoot/CS5489/exam-review/Midterm/items/MT047/question-en.tex}
\subsection*{中文题意}
\StudyTarget{WrittenResponse___20}
\input{\StudyRoot/CS5489/exam-review/Midterm/items/MT047/question-zh.tex}
\subsection*{English answer}
\StudyTarget{WrittenResponse__english-answer_9}
\input{\StudyRoot/CS5489/exam-review/Midterm/items/MT047/answer-en.tex}
\subsection*{中文对应答案}
\StudyTarget{WrittenResponse___21}
\input{\StudyRoot/CS5489/exam-review/Midterm/items/MT047/answer-zh.tex}
\input{\StudyRoot/CS5489/exam-review/Midterm/items/MT047/concise-extra.tex}
\renewcommand{\StudyQuestionExtra}{\StudyTarget{WrittenResponse__mt024}}
\StudyQuestion{MT024}{损失与正则项：\ensuremath{\alpha}从哪里来？}{简答题 / Written response · 解释}
\input{\StudyRoot/CS5489/exam-review/Midterm/items/MT024/provenance.tex}
\subsection*{English question}
\StudyTarget{WrittenResponse__english-question_10}
\input{\StudyRoot/CS5489/exam-review/Midterm/items/MT024/question-en.tex}
\subsection*{中文题意}
\StudyTarget{WrittenResponse___22}
\input{\StudyRoot/CS5489/exam-review/Midterm/items/MT024/question-zh.tex}
\subsection*{English answer}
\StudyTarget{WrittenResponse__english-answer_10}
\input{\StudyRoot/CS5489/exam-review/Midterm/items/MT024/answer-en.tex}
\subsection*{中文对应答案}
\StudyTarget{WrittenResponse___23}
\input{\StudyRoot/CS5489/exam-review/Midterm/items/MT024/answer-zh.tex}
\input{\StudyRoot/CS5489/exam-review/Midterm/items/MT024/concise-extra.tex}
\renewcommand{\StudyQuestionExtra}{\StudyTarget{WrittenResponse__mt012}\StudyTarget{WrittenResponse__mt012-l1}}
\StudyQuestion{MT012}{L1逻辑回归：写目标并画边界}{简答题 / Written response · 列式、作图、解释}
\input{\StudyRoot/CS5489/exam-review/Midterm/items/MT012/provenance.tex}
\subsection*{English question}
\StudyTarget{WrittenResponse__english-question_11}
\input{\StudyRoot/CS5489/exam-review/Midterm/items/MT012/question-en.tex}
\subsection*{中文题意}
\StudyTarget{WrittenResponse___24}
\input{\StudyRoot/CS5489/exam-review/Midterm/items/MT012/question-zh.tex}
\subsection*{English answer}
\StudyTarget{WrittenResponse__english-answer_11}
\input{\StudyRoot/CS5489/exam-review/Midterm/items/MT012/answer-en.tex}
\subsection*{中文对应答案}
\StudyTarget{WrittenResponse___25}
\input{\StudyRoot/CS5489/exam-review/Midterm/items/MT012/answer-zh.tex}
\input{\StudyRoot/CS5489/exam-review/Midterm/items/MT012/concise-extra.tex}
\Needspace{7\baselineskip}\phantomsection\addcontentsline{toc}{chapter}{SVM与核方法}
\renewcommand{\StudyQuestionExtra}{\StudyTarget{WrittenResponse__mt048}\StudyTarget{WrittenResponse__mt048-kernel-trick}}
\StudyQuestion{MT048}{用一个例子讲清Kernel Trick}{简答题 / Written response · 解释、举例}
\input{\StudyRoot/CS5489/exam-review/Midterm/items/MT048/provenance.tex}
\subsection*{English question}
\StudyTarget{WrittenResponse__english-question_12}
\input{\StudyRoot/CS5489/exam-review/Midterm/items/MT048/question-en.tex}
\subsection*{中文题意}
\StudyTarget{WrittenResponse___26}
\input{\StudyRoot/CS5489/exam-review/Midterm/items/MT048/question-zh.tex}
\subsection*{English answer}
\StudyTarget{WrittenResponse__english-answer_12}
\input{\StudyRoot/CS5489/exam-review/Midterm/items/MT048/answer-en.tex}
\subsection*{中文对应答案}
\StudyTarget{WrittenResponse___27}
\input{\StudyRoot/CS5489/exam-review/Midterm/items/MT048/answer-zh.tex}
\input{\StudyRoot/CS5489/exam-review/Midterm/items/MT048/concise-extra.tex}
\renewcommand{\StudyQuestionExtra}{\StudyTarget{WrittenResponse__mt060}\StudyTarget{WrittenResponse__mt060-svmcc}}
\StudyQuestion{MT060}{二次核SVM：大C与小C的边界}{简答题 / Written response · 作图、比较、解释}
\input{\StudyRoot/CS5489/exam-review/Midterm/items/MT060/provenance.tex}
\subsection*{English question}
\StudyTarget{WrittenResponse__english-question_13}
\input{\StudyRoot/CS5489/exam-review/Midterm/items/MT060/question-en.tex}
\subsection*{中文题意}
\StudyTarget{WrittenResponse___28}
\input{\StudyRoot/CS5489/exam-review/Midterm/items/MT060/question-zh.tex}
\subsection*{English answer}
\StudyTarget{WrittenResponse__english-answer_13}
\input{\StudyRoot/CS5489/exam-review/Midterm/items/MT060/answer-en.tex}
\subsection*{中文对应答案}
\StudyTarget{WrittenResponse___29}
\input{\StudyRoot/CS5489/exam-review/Midterm/items/MT060/answer-zh.tex}
\input{\StudyRoot/CS5489/exam-review/Midterm/items/MT060/concise-extra.tex}
\renewcommand{\StudyQuestionExtra}{\StudyTarget{WrittenResponse__mt010}\StudyTarget{WrittenResponse__mt010-gaussian-bayessvm}}
\StudyQuestion{MT010}{同一张分布图，Gaussian Bayes与二次核SVM如何判断？}{简答题 / Written response · 读图或公式、比较}
\input{\StudyRoot/CS5489/exam-review/Midterm/items/MT010/provenance.tex}
\subsection*{English question}
\StudyTarget{WrittenResponse__english-question_14}
\input{\StudyRoot/CS5489/exam-review/Midterm/items/MT010/question-en.tex}
\subsection*{中文题意}
\StudyTarget{WrittenResponse___30}
\input{\StudyRoot/CS5489/exam-review/Midterm/items/MT010/question-zh.tex}
\subsection*{English answer}
\StudyTarget{WrittenResponse__english-answer_14}
\input{\StudyRoot/CS5489/exam-review/Midterm/items/MT010/answer-en.tex}
\subsection*{中文对应答案}
\StudyTarget{WrittenResponse___31}
\input{\StudyRoot/CS5489/exam-review/Midterm/items/MT010/answer-zh.tex}
\input{\StudyRoot/CS5489/exam-review/Midterm/items/MT010/concise-extra.tex}
\renewcommand{\StudyQuestionExtra}{\StudyTarget{WrittenResponse__mt009}\StudyTarget{WrittenResponse__mt009-svm}}
\StudyQuestion{MT009}{十万维特征、两千样本：SVM怎样提速？}{简答题 / Written response · 提出办法、解释}
\input{\StudyRoot/CS5489/exam-review/Midterm/items/MT009/provenance.tex}
\subsection*{English question}
\StudyTarget{WrittenResponse__english-question_15}
\input{\StudyRoot/CS5489/exam-review/Midterm/items/MT009/question-en.tex}
\subsection*{中文题意}
\StudyTarget{WrittenResponse___32}
\input{\StudyRoot/CS5489/exam-review/Midterm/items/MT009/question-zh.tex}
\subsection*{English answer}
\StudyTarget{WrittenResponse__english-answer_15}
\input{\StudyRoot/CS5489/exam-review/Midterm/items/MT009/answer-en.tex}
\subsection*{中文对应答案}
\StudyTarget{WrittenResponse___33}
\input{\StudyRoot/CS5489/exam-review/Midterm/items/MT009/answer-zh.tex}
\input{\StudyRoot/CS5489/exam-review/Midterm/items/MT009/concise-extra.tex}
\renewcommand{\StudyQuestionExtra}{\StudyTarget{WrittenResponse__mt061}\StudyTarget{WrittenResponse__mt061-svmrbf-svm}}
\StudyQuestion{MT061}{预测内存：线性SVM与RBF SVM}{简答题 / Written response · 排序、解释}
\input{\StudyRoot/CS5489/exam-review/Midterm/items/MT061/provenance.tex}
\subsection*{English question}
\StudyTarget{WrittenResponse__english-question_16}
\input{\StudyRoot/CS5489/exam-review/Midterm/items/MT061/question-en.tex}
\subsection*{中文题意}
\StudyTarget{WrittenResponse___34}
\input{\StudyRoot/CS5489/exam-review/Midterm/items/MT061/question-zh.tex}
\subsection*{English answer}
\StudyTarget{WrittenResponse__english-answer_16}
\input{\StudyRoot/CS5489/exam-review/Midterm/items/MT061/answer-en.tex}
\subsection*{中文对应答案}
\StudyTarget{WrittenResponse___35}
\input{\StudyRoot/CS5489/exam-review/Midterm/items/MT061/answer-zh.tex}
\input{\StudyRoot/CS5489/exam-review/Midterm/items/MT061/concise-extra.tex}
\Needspace{7\baselineskip}\phantomsection\addcontentsline{toc}{chapter}{回归、特征选择与残差损失}
\renewcommand{\StudyQuestionExtra}{\StudyTarget{WrittenResponse__mt025}}
\StudyQuestion{MT025}{绝对残差损失对异常值有什么影响？}{简答题 / Written response · 解释}
\input{\StudyRoot/CS5489/exam-review/Midterm/items/MT025/provenance.tex}
\subsection*{English question}
\StudyTarget{WrittenResponse__english-question_17}
\input{\StudyRoot/CS5489/exam-review/Midterm/items/MT025/question-en.tex}
\subsection*{中文题意}
\StudyTarget{WrittenResponse___36}
\input{\StudyRoot/CS5489/exam-review/Midterm/items/MT025/question-zh.tex}
\subsection*{English answer}
\StudyTarget{WrittenResponse__english-answer_17}
\input{\StudyRoot/CS5489/exam-review/Midterm/items/MT025/answer-en.tex}
\subsection*{中文对应答案}
\StudyTarget{WrittenResponse___37}
\input{\StudyRoot/CS5489/exam-review/Midterm/items/MT025/answer-zh.tex}
\input{\StudyRoot/CS5489/exam-review/Midterm/items/MT025/concise-extra.tex}
\renewcommand{\StudyQuestionExtra}{\StudyTarget{WrittenResponse__mt063}\StudyTarget{WrittenResponse__mt063-huber}}
\StudyQuestion{MT063}{Huber损失：画图并解释好处}{简答题 / Written response · 作图、解释}
\input{\StudyRoot/CS5489/exam-review/Midterm/items/MT063/provenance.tex}
\subsection*{English question}
\StudyTarget{WrittenResponse__english-question_18}
\input{\StudyRoot/CS5489/exam-review/Midterm/items/MT063/question-en.tex}
\subsection*{中文题意}
\StudyTarget{WrittenResponse___38}
\input{\StudyRoot/CS5489/exam-review/Midterm/items/MT063/question-zh.tex}
\subsection*{English answer}
\StudyTarget{WrittenResponse__english-answer_18}
\input{\StudyRoot/CS5489/exam-review/Midterm/items/MT063/answer-en.tex}
\subsection*{中文对应答案}
\StudyTarget{WrittenResponse___39}
\input{\StudyRoot/CS5489/exam-review/Midterm/items/MT063/answer-zh.tex}
\input{\StudyRoot/CS5489/exam-review/Midterm/items/MT063/concise-extra.tex}
\renewcommand{\StudyQuestionExtra}{\StudyTarget{WrittenResponse__mt013}\StudyTarget{WrittenResponse__mt013-l1l2}}
\StudyQuestion{MT013}{残差L1＋权重L2，各自在做什么？}{简答题 / Written response · 解释}
\input{\StudyRoot/CS5489/exam-review/Midterm/items/MT013/provenance.tex}
\subsection*{English question}
\StudyTarget{WrittenResponse__english-question_19}
\input{\StudyRoot/CS5489/exam-review/Midterm/items/MT013/question-en.tex}
\subsection*{中文题意}
\StudyTarget{WrittenResponse___40}
\input{\StudyRoot/CS5489/exam-review/Midterm/items/MT013/question-zh.tex}
\subsection*{English answer}
\StudyTarget{WrittenResponse__english-answer_19}
\input{\StudyRoot/CS5489/exam-review/Midterm/items/MT013/answer-en.tex}
\subsection*{中文对应答案}
\StudyTarget{WrittenResponse___41}
\input{\StudyRoot/CS5489/exam-review/Midterm/items/MT013/answer-zh.tex}
\input{\StudyRoot/CS5489/exam-review/Midterm/items/MT013/concise-extra.tex}
\renewcommand{\StudyQuestionExtra}{\StudyTarget{WrittenResponse__mt050}}
\StudyQuestion{MT050}{缺车比车多更糟：设计非对称损失}{简答题 / Written response · 设计、作图、解释}
\input{\StudyRoot/CS5489/exam-review/Midterm/items/MT050/provenance.tex}
\subsection*{English question}
\StudyTarget{WrittenResponse__english-question_20}
\input{\StudyRoot/CS5489/exam-review/Midterm/items/MT050/question-en.tex}
\subsection*{中文题意}
\StudyTarget{WrittenResponse___42}
\input{\StudyRoot/CS5489/exam-review/Midterm/items/MT050/question-zh.tex}
\subsection*{English answer}
\StudyTarget{WrittenResponse__english-answer_20}
\input{\StudyRoot/CS5489/exam-review/Midterm/items/MT050/answer-en.tex}
\subsection*{中文对应答案}
\StudyTarget{WrittenResponse___43}
\input{\StudyRoot/CS5489/exam-review/Midterm/items/MT050/answer-zh.tex}
\input{\StudyRoot/CS5489/exam-review/Midterm/items/MT050/concise-extra.tex}
\renewcommand{\StudyQuestionExtra}{\StudyTarget{WrittenResponse__mt037}}
\StudyQuestion{MT037}{线性回归为何与均值预测器表现相近？}{简答题 / Written response · 提出办法、解释}
\input{\StudyRoot/CS5489/exam-review/Midterm/items/MT037/provenance.tex}
\subsection*{English question}
\StudyTarget{WrittenResponse__english-question_21}
\input{\StudyRoot/CS5489/exam-review/Midterm/items/MT037/question-en.tex}
\subsection*{中文题意}
\StudyTarget{WrittenResponse___44}
\input{\StudyRoot/CS5489/exam-review/Midterm/items/MT037/question-zh.tex}
\subsection*{English answer}
\StudyTarget{WrittenResponse__english-answer_21}
\input{\StudyRoot/CS5489/exam-review/Midterm/items/MT037/answer-en.tex}
\subsection*{中文对应答案}
\StudyTarget{WrittenResponse___45}
\input{\StudyRoot/CS5489/exam-review/Midterm/items/MT037/answer-zh.tex}
\input{\StudyRoot/CS5489/exam-review/Midterm/items/MT037/concise-extra.tex}
\Needspace{7\baselineskip}\phantomsection\addcontentsline{toc}{chapter}{集成学习}
\renewcommand{\StudyQuestionExtra}{\StudyTarget{WrittenResponse__mt062}}
\StudyQuestion{MT062}{随机森林：树多了为什么仍有共同误差？}{简答题 / Written response · 读图或公式、提出办法}
\input{\StudyRoot/CS5489/exam-review/Midterm/items/MT062/provenance.tex}
\subsection*{English question}
\StudyTarget{WrittenResponse__english-question_22}
\input{\StudyRoot/CS5489/exam-review/Midterm/items/MT062/question-en.tex}
\subsection*{中文题意}
\StudyTarget{WrittenResponse___46}
\input{\StudyRoot/CS5489/exam-review/Midterm/items/MT062/question-zh.tex}
\subsection*{English answer}
\StudyTarget{WrittenResponse__english-answer_22}
\input{\StudyRoot/CS5489/exam-review/Midterm/items/MT062/answer-en.tex}
\subsection*{中文对应答案}
\StudyTarget{WrittenResponse___47}
\input{\StudyRoot/CS5489/exam-review/Midterm/items/MT062/answer-zh.tex}
\input{\StudyRoot/CS5489/exam-review/Midterm/items/MT062/concise-extra.tex}
\renewcommand{\StudyQuestionExtra}{\StudyTarget{WrittenResponse__mt038}\StudyTarget{WrittenResponse__mt038-100rfadaboost}}
\StudyQuestion{MT038}{100核部署：RF与AdaBoost如何并行预测？}{简答题 / Written response · 比较、提出办法}
\input{\StudyRoot/CS5489/exam-review/Midterm/items/MT038/provenance.tex}
\subsection*{English question}
\StudyTarget{WrittenResponse__english-question_23}
\input{\StudyRoot/CS5489/exam-review/Midterm/items/MT038/question-en.tex}
\subsection*{中文题意}
\StudyTarget{WrittenResponse___48}
\input{\StudyRoot/CS5489/exam-review/Midterm/items/MT038/question-zh.tex}
\subsection*{English answer}
\StudyTarget{WrittenResponse__english-answer_23}
\input{\StudyRoot/CS5489/exam-review/Midterm/items/MT038/answer-en.tex}
\subsection*{中文对应答案}
\StudyTarget{WrittenResponse___49}
\input{\StudyRoot/CS5489/exam-review/Midterm/items/MT038/answer-zh.tex}
\input{\StudyRoot/CS5489/exam-review/Midterm/items/MT038/concise-extra.tex}
\renewcommand{\StudyQuestionExtra}{\StudyTarget{WrittenResponse__mt049}\StudyTarget{WrittenResponse__mt049-adaboost}}
\StudyQuestion{MT049}{在线更新AdaBoost：速度与异常值}{简答题 / Written response · 提出办法、判断并说明}
\input{\StudyRoot/CS5489/exam-review/Midterm/items/MT049/provenance.tex}
\subsection*{English question}
\StudyTarget{WrittenResponse__english-question_24}
\input{\StudyRoot/CS5489/exam-review/Midterm/items/MT049/question-en.tex}
\subsection*{中文题意}
\StudyTarget{WrittenResponse___50}
\input{\StudyRoot/CS5489/exam-review/Midterm/items/MT049/question-zh.tex}
\subsection*{English answer}
\StudyTarget{WrittenResponse__english-answer_24}
\input{\StudyRoot/CS5489/exam-review/Midterm/items/MT049/answer-en.tex}
\subsection*{中文对应答案}
\StudyTarget{WrittenResponse___51}
\input{\StudyRoot/CS5489/exam-review/Midterm/items/MT049/answer-zh.tex}
\input{\StudyRoot/CS5489/exam-review/Midterm/items/MT049/concise-extra.tex}
\renewcommand{\StudyQuestionExtra}{\StudyTarget{WrittenResponse__mt064}\StudyTarget{WrittenResponse__mt064-baggingboostinghuber}}
\StudyQuestion{MT064}{综合题：Bagging/Boosting与Huber}{简答题 / Written response · 比较、作图、解释}
\input{\StudyRoot/CS5489/exam-review/Midterm/items/MT064/provenance.tex}
\subsection*{English question}
\StudyTarget{WrittenResponse__english-question_25}
\input{\StudyRoot/CS5489/exam-review/Midterm/items/MT064/question-en.tex}
\subsection*{中文题意}
\StudyTarget{WrittenResponse___52}
\input{\StudyRoot/CS5489/exam-review/Midterm/items/MT064/question-zh.tex}
\subsection*{English answer}
\StudyTarget{WrittenResponse__english-answer_25}
\input{\StudyRoot/CS5489/exam-review/Midterm/items/MT064/answer-en.tex}
\subsection*{中文对应答案}
\StudyTarget{WrittenResponse___53}
\input{\StudyRoot/CS5489/exam-review/Midterm/items/MT064/answer-zh.tex}
\input{\StudyRoot/CS5489/exam-review/Midterm/items/MT064/concise-extra.tex}
\Needspace{7\baselineskip}\phantomsection\addcontentsline{toc}{chapter}{分类损失图与错误代价}
\renewcommand{\StudyQuestionExtra}{\StudyTarget{WrittenResponse__mt008}}
\StudyQuestion{MT008}{平方形分类损失为什么惩罚“太正确”？}{简答题 / Written response · 读图或公式、判断并说明}
\input{\StudyRoot/CS5489/exam-review/Midterm/items/MT008/provenance.tex}
\subsection*{English question}
\StudyTarget{WrittenResponse__english-question_26}
\input{\StudyRoot/CS5489/exam-review/Midterm/items/MT008/question-en.tex}
\subsection*{中文题意}
\StudyTarget{WrittenResponse___54}
\input{\StudyRoot/CS5489/exam-review/Midterm/items/MT008/question-zh.tex}
\subsection*{English answer}
\StudyTarget{WrittenResponse__english-answer_26}
\input{\StudyRoot/CS5489/exam-review/Midterm/items/MT008/answer-en.tex}
\subsection*{中文对应答案}
\StudyTarget{WrittenResponse___55}
\input{\StudyRoot/CS5489/exam-review/Midterm/items/MT008/answer-zh.tex}
\input{\StudyRoot/CS5489/exam-review/Midterm/items/MT008/concise-extra.tex}
\renewcommand{\StudyQuestionExtra}{\StudyTarget{WrittenResponse__mt026}\StudyTarget{WrittenResponse__mt026-vz1}}
\StudyQuestion{MT026}{V形损失：最优点为何在z=1？}{简答题 / Written response · 读图或公式、判断并说明}
\input{\StudyRoot/CS5489/exam-review/Midterm/items/MT026/provenance.tex}
\subsection*{English question}
\StudyTarget{WrittenResponse__english-question_27}
\input{\StudyRoot/CS5489/exam-review/Midterm/items/MT026/question-en.tex}
\subsection*{中文题意}
\StudyTarget{WrittenResponse___56}
\input{\StudyRoot/CS5489/exam-review/Midterm/items/MT026/question-zh.tex}
\subsection*{English answer}
\StudyTarget{WrittenResponse__english-answer_27}
\input{\StudyRoot/CS5489/exam-review/Midterm/items/MT026/answer-en.tex}
\subsection*{中文对应答案}
\StudyTarget{WrittenResponse___57}
\input{\StudyRoot/CS5489/exam-review/Midterm/items/MT026/answer-zh.tex}
\input{\StudyRoot/CS5489/exam-review/Midterm/items/MT026/concise-extra.tex}
\renewcommand{\StudyQuestionExtra}{\StudyTarget{WrittenResponse__mt051}}
\StudyQuestion{MT051}{负间隔也零损失：分类器会放过什么？}{简答题 / Written response · 读图或公式、判断并说明、作图}
\input{\StudyRoot/CS5489/exam-review/Midterm/items/MT051/provenance.tex}
\subsection*{English question}
\StudyTarget{WrittenResponse__english-question_28}
\input{\StudyRoot/CS5489/exam-review/Midterm/items/MT051/question-en.tex}
\subsection*{中文题意}
\StudyTarget{WrittenResponse___58}
\input{\StudyRoot/CS5489/exam-review/Midterm/items/MT051/question-zh.tex}
\subsection*{English answer}
\StudyTarget{WrittenResponse__english-answer_28}
\input{\StudyRoot/CS5489/exam-review/Midterm/items/MT051/answer-en.tex}
\subsection*{中文对应答案}
\StudyTarget{WrittenResponse___59}
\input{\StudyRoot/CS5489/exam-review/Midterm/items/MT051/answer-zh.tex}
\input{\StudyRoot/CS5489/exam-review/Midterm/items/MT051/concise-extra.tex}
\renewcommand{\StudyQuestionExtra}{\StudyTarget{WrittenResponse__mt039}}
\StudyQuestion{MT039}{饱和的非凸分类损失}{简答题 / Written response · 读图或公式}
\input{\StudyRoot/CS5489/exam-review/Midterm/items/MT039/provenance.tex}
\subsection*{English question}
\StudyTarget{WrittenResponse__english-question_29}
\input{\StudyRoot/CS5489/exam-review/Midterm/items/MT039/question-en.tex}
\subsection*{中文题意}
\StudyTarget{WrittenResponse___60}
\input{\StudyRoot/CS5489/exam-review/Midterm/items/MT039/question-zh.tex}
\subsection*{English answer}
\StudyTarget{WrittenResponse__english-answer_29}
\input{\StudyRoot/CS5489/exam-review/Midterm/items/MT039/answer-en.tex}
\subsection*{中文对应答案}
\StudyTarget{WrittenResponse___61}
\input{\StudyRoot/CS5489/exam-review/Midterm/items/MT039/answer-zh.tex}
\input{\StudyRoot/CS5489/exam-review/Midterm/items/MT039/concise-extra.tex}
\renewcommand{\StudyQuestionExtra}{\StudyTarget{WrittenResponse__mt065}\StudyTarget{WrittenResponse__mt065-logistichingeexponential}}
\StudyQuestion{MT065}{比较Logistic、Hinge和Exponential损失}{简答题 / Written response · 比较、排序}
\input{\StudyRoot/CS5489/exam-review/Midterm/items/MT065/provenance.tex}
\subsection*{English question}
\StudyTarget{WrittenResponse__english-question_30}
\input{\StudyRoot/CS5489/exam-review/Midterm/items/MT065/question-en.tex}
\subsection*{中文题意}
\StudyTarget{WrittenResponse___62}
\input{\StudyRoot/CS5489/exam-review/Midterm/items/MT065/question-zh.tex}
\subsection*{English answer}
\StudyTarget{WrittenResponse__english-answer_30}
\input{\StudyRoot/CS5489/exam-review/Midterm/items/MT065/answer-en.tex}
\subsection*{中文对应答案}
\StudyTarget{WrittenResponse___63}
\input{\StudyRoot/CS5489/exam-review/Midterm/items/MT065/answer-zh.tex}
\input{\StudyRoot/CS5489/exam-review/Midterm/items/MT065/concise-extra.tex}
\renewcommand{\StudyQuestionExtra}{\StudyTarget{WrittenResponse__mt007}}
\StudyQuestion{MT007}{筛查任务：漏诊比误报更贵}{简答题 / Written response · 提出办法、解释}
\input{\StudyRoot/CS5489/exam-review/Midterm/items/MT007/provenance.tex}
\subsection*{English question}
\StudyTarget{WrittenResponse__english-question_31}
\input{\StudyRoot/CS5489/exam-review/Midterm/items/MT007/question-en.tex}
\subsection*{中文题意}
\StudyTarget{WrittenResponse___64}
\input{\StudyRoot/CS5489/exam-review/Midterm/items/MT007/question-zh.tex}
\subsection*{English answer}
\StudyTarget{WrittenResponse__english-answer_31}
\input{\StudyRoot/CS5489/exam-review/Midterm/items/MT007/answer-en.tex}
\subsection*{中文对应答案}
\StudyTarget{WrittenResponse___65}
\input{\StudyRoot/CS5489/exam-review/Midterm/items/MT007/answer-zh.tex}
\input{\StudyRoot/CS5489/exam-review/Midterm/items/MT007/concise-extra.tex}
\renewcommand{\StudyQuestionExtra}{\StudyTarget{WrittenResponse__mt023}\StudyTarget{WrittenResponse__mt023-100050}}
\StudyQuestion{MT023}{1000名受试者中只有50名阳性}{简答题 / Written response · 诊断问题、提出办法}
\input{\StudyRoot/CS5489/exam-review/Midterm/items/MT023/provenance.tex}
\subsection*{English question}
\StudyTarget{WrittenResponse__english-question_32}
\input{\StudyRoot/CS5489/exam-review/Midterm/items/MT023/question-en.tex}
\subsection*{中文题意}
\StudyTarget{WrittenResponse___66}
\input{\StudyRoot/CS5489/exam-review/Midterm/items/MT023/question-zh.tex}
\subsection*{English answer}
\StudyTarget{WrittenResponse__english-answer_32}
\input{\StudyRoot/CS5489/exam-review/Midterm/items/MT023/answer-en.tex}
\subsection*{中文对应答案}
\StudyTarget{WrittenResponse___67}
\input{\StudyRoot/CS5489/exam-review/Midterm/items/MT023/answer-zh.tex}
\input{\StudyRoot/CS5489/exam-review/Midterm/items/MT023/concise-extra.tex}
\renewcommand{\StudyQuestionExtra}{\StudyTarget{WrittenResponse__mt035}}
\StudyQuestion{MT035}{正常邮件被大量拦截怎么办？}{简答题 / Written response · 提出办法、解释}
\input{\StudyRoot/CS5489/exam-review/Midterm/items/MT035/provenance.tex}
\subsection*{English question}
\StudyTarget{WrittenResponse__english-question_33}
\input{\StudyRoot/CS5489/exam-review/Midterm/items/MT035/question-en.tex}
\subsection*{中文题意}
\StudyTarget{WrittenResponse___68}
\input{\StudyRoot/CS5489/exam-review/Midterm/items/MT035/question-zh.tex}
\subsection*{English answer}
\StudyTarget{WrittenResponse__english-answer_33}
\input{\StudyRoot/CS5489/exam-review/Midterm/items/MT035/answer-en.tex}
\subsection*{中文对应答案}
\StudyTarget{WrittenResponse___69}
\input{\StudyRoot/CS5489/exam-review/Midterm/items/MT035/answer-zh.tex}
\input{\StudyRoot/CS5489/exam-review/Midterm/items/MT035/concise-extra.tex}
\Needspace{7\baselineskip}\phantomsection\addcontentsline{toc}{chapter}{神经网络与链式法则}
\renewcommand{\StudyQuestionExtra}{\StudyTarget{WrittenResponse__mt067}}
\StudyQuestion{MT067}{链式求导与两层网络反传}{简答题 / Written response · 计算推导}
\input{\StudyRoot/CS5489/exam-review/Midterm/items/MT067/provenance.tex}
\subsection*{English question}
\StudyTarget{WrittenResponse__english-question_34}
\input{\StudyRoot/CS5489/exam-review/Midterm/items/MT067/question-en.tex}
\subsection*{中文题意}
\StudyTarget{WrittenResponse___70}
\input{\StudyRoot/CS5489/exam-review/Midterm/items/MT067/question-zh.tex}
\subsection*{English answer}
\StudyTarget{WrittenResponse__english-answer_34}
\input{\StudyRoot/CS5489/exam-review/Midterm/items/MT067/answer-en.tex}
\subsection*{中文对应答案}
\StudyTarget{WrittenResponse___71}
\input{\StudyRoot/CS5489/exam-review/Midterm/items/MT067/answer-zh.tex}
\input{\StudyRoot/CS5489/exam-review/Midterm/items/MT067/concise-extra.tex}
\renewcommand{\StudyQuestionExtra}{\StudyTarget{WrittenResponse__mt066}}
\StudyQuestion{MT066}{梯度消失：为什么深度和激活函数有关？}{简答题 / Written response · 解释、计算推导}
\input{\StudyRoot/CS5489/exam-review/Midterm/items/MT066/provenance.tex}
\subsection*{English question}
\StudyTarget{WrittenResponse__english-question_35}
\input{\StudyRoot/CS5489/exam-review/Midterm/items/MT066/question-en.tex}
\subsection*{中文题意}
\StudyTarget{WrittenResponse___72}
\input{\StudyRoot/CS5489/exam-review/Midterm/items/MT066/question-zh.tex}
\subsection*{English answer}
\StudyTarget{WrittenResponse__english-answer_35}
\input{\StudyRoot/CS5489/exam-review/Midterm/items/MT066/answer-en.tex}
\subsection*{中文对应答案}
\StudyTarget{WrittenResponse___73}
\input{\StudyRoot/CS5489/exam-review/Midterm/items/MT066/answer-zh.tex}
\input{\StudyRoot/CS5489/exam-review/Midterm/items/MT066/concise-extra.tex}
\Needspace{7\baselineskip}\phantomsection\addcontentsline{toc}{chapter}{历史范围拓展：Gaussian Process}
\renewcommand{\StudyQuestionExtra}{\StudyTarget{WrittenResponse__mt036}\StudyTarget{WrittenResponse__mt036-gprf}}
\StudyQuestion{MT036}{历史拓展：音乐地理回归，GP还是RF？}{简答题 / Written response · 比较}
\input{\StudyRoot/CS5489/exam-review/Midterm/items/MT036/provenance.tex}
\subsection*{English question}
\StudyTarget{WrittenResponse__english-question_36}
\input{\StudyRoot/CS5489/exam-review/Midterm/items/MT036/question-en.tex}
\subsection*{中文题意}
\StudyTarget{WrittenResponse___74}
\input{\StudyRoot/CS5489/exam-review/Midterm/items/MT036/question-zh.tex}
\subsection*{English answer}
\StudyTarget{WrittenResponse__english-answer_36}
\input{\StudyRoot/CS5489/exam-review/Midterm/items/MT036/answer-en.tex}
\subsection*{中文对应答案}
\StudyTarget{WrittenResponse___75}
\input{\StudyRoot/CS5489/exam-review/Midterm/items/MT036/answer-zh.tex}
\input{\StudyRoot/CS5489/exam-review/Midterm/items/MT036/concise-extra.tex}
\clearpage\StudyHeading{1}{原卷索引 / Original-paper index}
按原卷题号查找本册位置；不同年份的重复题指向同一MT编号。2021B*保留题纸与配套答案的年份冲突。
\StudyTarget{WrittenResponse__original-paper-index}
\subsection*{2020B期中Quiz}
\StudyTarget{WrittenResponse__2020bquiz}
\begin{longtable}{@{}p{.2\linewidth}p{.56\linewidth}p{.18\linewidth}@{}}
\toprule 原题 & 本册题号与页码 & 分值 \\ \midrule\endhead
Q7 & \StudyRef{WrittenResponse}{mt007}{MT007} & 10 \\
Q8 & \StudyRef{WrittenResponse}{mt008}{MT008} & 10 \\
Q9 & \StudyRef{WrittenResponse}{mt009}{MT009} & 10 \\
Q10 & \StudyRef{WrittenResponse}{mt010}{MT010} & 10 \\
Q11 & \StudyRef{WrittenResponse}{mt011}{MT011} & 10 \\
Q12 & \StudyRef{WrittenResponse}{mt012}{MT012} & 10 \\
Q13 & \StudyRef{WrittenResponse}{mt013}{MT013} & 10 \\
\bottomrule\end{longtable}
\subsection*{2021A期中}
\StudyTarget{WrittenResponse__2021a}
\begin{longtable}{@{}p{.2\linewidth}p{.56\linewidth}p{.18\linewidth}@{}}
\toprule 原题 & 本册题号与页码 & 分值 \\ \midrule\endhead
Q7 & \StudyRef{WrittenResponse}{mt020}{MT020} & 10 \\
Q8 & \StudyRef{WrittenResponse}{mt021}{MT021} & 10 \\
Q9 & \StudyRef{WrittenResponse}{mt022}{MT022} & 10 \\
Q10 & \StudyRef{WrittenResponse}{mt023}{MT023} & 10 \\
Q11 & \StudyRef{WrittenResponse}{mt024}{MT024} & 10 \\
Q12 & \StudyRef{WrittenResponse}{mt025}{MT025} & 10 \\
Q13 & \StudyRef{WrittenResponse}{mt026}{MT026} & 10 \\
\bottomrule\end{longtable}
\subsection*{2021B*期中}
\StudyTarget{WrittenResponse__2021b}
\begin{longtable}{@{}p{.2\linewidth}p{.56\linewidth}p{.18\linewidth}@{}}
\toprule 原题 & 本册题号与页码 & 分值 \\ \midrule\endhead
Q7 & \StudyRef{WrittenResponse}{mt033}{MT033} & 10 \\
Q8 & \StudyRef{WrittenResponse}{mt034}{MT034} & 10 \\
Q9 & \StudyRef{WrittenResponse}{mt035}{MT035} & 10 \\
Q10 & \StudyRef{WrittenResponse}{mt036}{MT036} & 10 \\
Q11 & \StudyRef{WrittenResponse}{mt037}{MT037} & 10 \\
Q12 & \StudyRef{WrittenResponse}{mt038}{MT038} & 10 \\
Q13 & \StudyRef{WrittenResponse}{mt039}{MT039} & 10 \\
\bottomrule\end{longtable}
\subsection*{2023A期中}
\StudyTarget{WrittenResponse__2023a}
\begin{longtable}{@{}p{.2\linewidth}p{.56\linewidth}p{.18\linewidth}@{}}
\toprule 原题 & 本册题号与页码 & 分值 \\ \midrule\endhead
Q7 & \StudyRef{WrittenResponse}{mt046}{MT046} & 10 \\
Q8 & \StudyRef{WrittenResponse}{mt047}{MT047} & 10 \\
Q9 & \StudyRef{WrittenResponse}{mt048}{MT048} & 10 \\
Q10 & \StudyRef{WrittenResponse}{mt023}{MT023} & 10 \\
Q11 & \StudyRef{WrittenResponse}{mt049}{MT049} & 10 \\
Q12 & \StudyRef{WrittenResponse}{mt050}{MT050} & 10 \\
Q13 & \StudyRef{WrittenResponse}{mt051}{MT051} & 10 \\
\bottomrule\end{longtable}
\subsection*{2023B期中}
\StudyTarget{WrittenResponse__2023b}
\begin{longtable}{@{}p{.2\linewidth}p{.56\linewidth}p{.18\linewidth}@{}}
\toprule 原题 & 本册题号与页码 & 分值 \\ \midrule\endhead
Q7 & \StudyRef{WrittenResponse}{mt058}{MT058} & 10 \\
Q8 & \StudyRef{WrittenResponse}{mt059}{MT059} & 10 \\
Q9 & \StudyRef{WrittenResponse}{mt060}{MT060} & 10 \\
Q10 & \StudyRef{WrittenResponse}{mt061}{MT061} & 10 \\
Q11 & \StudyRef{WrittenResponse}{mt062}{MT062} & 10 \\
Q12 & \StudyRef{WrittenResponse}{mt063}{MT063} & 10 \\
Q13 & \StudyRef{WrittenResponse}{mt007}{MT007} & 10 \\
\bottomrule\end{longtable}
\subsection*{2025A期中}
\StudyTarget{WrittenResponse__2025a}
\begin{longtable}{@{}p{.2\linewidth}p{.56\linewidth}p{.18\linewidth}@{}}
\toprule 原题 & 本册题号与页码 & 分值 \\ \midrule\endhead
Q7 & \StudyRef{WrittenResponse}{mt064}{MT064} & 10 \\
Q8 & \StudyRef{WrittenResponse}{mt047}{MT047} & 10 \\
Q9 & \StudyRef{WrittenResponse}{mt065}{MT065} & 10 \\
Q10 & \StudyRef{WrittenResponse}{mt066}{MT066} & 10 \\
Q11 & \StudyRef{WrittenResponse}{mt046}{MT046} & 10 \\
Q12 & \StudyRef{WrittenResponse}{mt050}{MT050} & 10 \\
Q13 & \StudyRef{WrittenResponse}{mt067}{MT067} & 10 \\
\bottomrule\end{longtable}
\end{document}
